\section{The common differentiated activity bound}
\label{sec:activity-jet-assembly}
We assemble the estimates on the actual four-part alphabet:
connected seed packets, forced activities, exponentiated
logarithmic activities, and warnings. All factors use the same
sector, fine representation, tags, Gaussian law and compensated
variation. Besides the hypotheses of the preceding local estimates,
retain the deterministic safety and plateau inequalities of
Section~\ref{sec:gaussian-warnings}. Their construction is an
incoming requirement at this stage. The argument below explains
both their derivative use and the common price they yield.

\subsection{Exponentiating the logarithmic activities}
If $f\in\mathcal B_j(w,d;y_0)$ with $0\le w\le1$, then
\begin{equation}\label{eq:activity-exp-jets}
 e^f-1\in\mathcal B_j(C_jw,d;y_0)
\end{equation}
for a constant depending only on $j$. For the value, the mean
value theorem gives $|e^{f(y_0)}-1|\le e w$. For a positive
derivative order $k\le j$, use the partition formula for the
exponential. A partition with $m\ge1$ blocks contributes at
most $e w^m d^k\le e w d^k$. There are at most $k!$ partitions,
so $C_j=e\max(1,j!)$ suffices. This retains the small amplitude
and does not require that the parameter variation stay within
the small-amplitude bound on a whole neighbourhood.

Fix $-1/2<t<0$. The inherited logarithmic prices of
Section~\ref{sec:inherited-activity-prices}, the three short
prices of Section~\ref{sec:remote-action-jets}, and the correction
chain prices of Section~\ref{sec:itinerary-prices} have a common
rooted bound $J(x)$ of upper exponent $t$, after an arbitrarily
small strict exponent loss. Their supports are nonempty. Since
the support exponent and extra-cost weight are nonnegative,
each individual price is at most $J(x)$: choose one site in its
support and retain that summand in the rooted bound. Eventually
$J(x)\le1$. Thus \eqref{eq:activity-exp-jets} applies to every
logarithmic label with the same $C_j$ and the same threshold,
even when the alphabet is countably infinite. Multiplying all
its prices by $C_j$ preserves their upper exponent.

The forced family needs no exponentiation. By enlarging a
nonnegative function $r(x)\ge1$ of upper exponent zero, both
families have jets with derivative scale
\begin{equation}\label{eq:activity-common-scale}
 r(x)\{1+c_B(S_i)+d_i\}.
\end{equation}
Here $S_i$ is the full assigned support and $d_i$ its old-cover
or path-length charge; for a correction itinerary it is the
correction order. The short terms have $d_i=0$. All earlier
scale bounds are dominated by this expression.

\subsection{Warnings vanish as germs at a reduced threshold}
Let $\xi>0$ be the safety margin and use the smaller threshold
$\xi/2$ in the Gaussian event. Suppose every coordinate within
distance $2B$ of a warning anchor satisfies
\[
 |z_i|\le(\xi/2)\tau\sqrt\beta.
\]
The deterministic safety estimate at margin $\xi$ implies that
the warning factor is zero on a neighbourhood of the base
parameter. We check the different dependencies, since a value
equal to zero by itself would not imply vanishing derivatives.

A warning in its mandatory seed neighbourhood is identically
zero by its frozen indicator. Outside that neighbourhood, the
safety estimate makes every relevant fine plaquette reading
strictly smaller than $\tau/4$, and every predictor chord
strictly smaller than $W\tau$. The fine mask, site likelihood,
bond likelihood and contraction test are evaluated at the held
fine configuration $U_*$. Their values are consequently constant
under the variation and equal to one. The contraction statement
uses the fine plaquette estimates at all sites in its read set,
which fits in the same $2B$ event window. The predictor chord
depends continuously on the varied coarse field, with its fine
predictor held fixed. Its strict inequality persists on a
neighbourhood, so its cutoff also remains one there.

For a fine-bond cutoff, the reduced threshold bounds the base
increment by $(C+q\xi/2)\tau$. The safety constants satisfy
$C+q\xi/2<a$, where $a\tau$ is its plateau radius.
The compensated increment is continuous, so this strict
inequality also persists locally. Its cutoff-minus-one is thus
a zero germ. Intersecting the finitely many neighbourhoods
needed by a given factor proves the assertion. In particular,
every nonzero jet of a warning is supported on the reduced-margin
excursion event; no assertion about differentiation of an
indicator of that event is needed.

The warning amplitude is bounded by the fixed constant
\[
 T=1+\#\{\text{frame tags}\}+\#\{\text{predictor tags}\}.
\]
Indeed the held likelihoods are bounded by their tag counts,
all masks lie in $[0,1]$, and the warning is a protected
difference from one. Its derivative scale has the form
\eqref{eq:activity-common-scale} with $d_i=0$.
For a varying noise cutoff, normalize its chord by $W\tau$;
the positive derivatives cost the coarse scale divided by
$\tau$, of upper exponent zero. For a bond cutoff use the
corresponding normalized compensated increment. A cutoff
outside its closed support is a zero germ; on the closed support
the smooth scalar cutoff has bounded derivatives. Multiplying
over its protective read set adds at most a function of upper
exponent zero times its support cover. This proves the asserted
scale in all cases, including boundary points.

For every $t<0$, the Gaussian estimate at margin $\xi/2$ gives,
after increasing the threshold in $x$,
\begin{equation}\label{eq:activity-warning-square}
 \mathbb P(\text{all $m$ reduced-margin warnings})
       \le(h^{2t})^m.
\end{equation}
This follows from the faster-than-every-power bound already
proved, applied at an exponent strictly below $2t$.
Set $\rho=h^t$. The Gaussian moment argument consumes the
square root of \eqref{eq:activity-warning-square}, leaving
the price $T\rho$ per distinct warning. Its site sum has upper
exponent $t$: warning addresses have fixed density, each support
lies in a fixed-multiple-$B$ ball, its cover is bounded by a fixed
constant, and the number of such supports through a site is
polynomial in $B$.

\subsection{Packets and the product estimate}
The packet jets have price
\[
 p_A=\exp\{((C+3j)d^5-\beta\mu\tau^2)n(A)\}
\]
and derivative scale one by \eqref{eq:packet-jet-price}.
The same energy gain pays their root counts, support exponents
and any prescribed fixed seed marks by the absorption argument
in Section~\ref{sec:packet-estimates}. Thus their available
weighted site price is smaller than every fixed negative power
of $h$. A disconnected packet has factor and price zero.
All thresholds are chosen after the fixed derivative order,
halo, support and length reserves.

Let $\mathcal P$ be a finite family of distinct labels of the
four types whose assigned supports are healthy for the base coarse
field: no coarse flag at the threshold $\operatorname{temp}(x)$
lies in any $S_i$. In particular, each packet's full enlarged
support must have this property to apply its packet jet estimate.
This is a separate hypothesis on the tested family; the sector's
manifest flag identity alone does not imply it. Put
\[
 s=\sum_{i\in\mathcal P}\{c_B(S_i)+d_i\}.
\]
Every nonzero factor has nonempty support, hence $c_B(S_i)\ge1$.
The derivative scales in \eqref{eq:activity-common-scale}
therefore sum to at most $2r(x)s$, also covering the unit packet
scales. Let $a_i$ denote the packet, forced or exponentiated-log
price, and put $a_i=T$ for a warning. The product rule and the
zero-germ property give exactly
\[
 \prod_{i\in\mathcal P}f_i
 \in\mathcal B_j\!\left(
 \mathbf1_{\mathcal C}(z)\prod_{i\in\mathcal P}a_i,
 2r(x)s;y_0\right),
\]
where $\mathcal C$ is the joint reduced-margin warning event.
This verifies \eqref{eq:moment-activity-premise} for all species
on the same variation. Applying the Gaussian differentiation
and moment estimates yields
\begin{equation}\label{eq:activity-integrated-jet}
 \left\|D^k\int\prod_{i\in\mathcal P}f_i\right\|
 \le A_j(x)(1+s)^{m_j}
             \prod_{i\in\mathcal P}\widetilde a_i,
 \qquad 0\le k\le j,
\end{equation}
with $A_j$ of upper exponent zero and
$\widetilde a_i=a_i$ except that a warning has price $T\rho$.
The integral includes the Gaussian dependence on the varied
field and the frozen Haar variables, as in that argument.

The combined optional site sum has upper exponent $t$, and
the packet site sum is negligible. For a desired output exponent
$t_{\rm out}\in(-1/2,0)$, choose $t\in(-1/2,t_{\rm out})$. Fixed
upper-exponent-zero factors and the fixed polynomial in $1+s$
are then paid by the strict exponent slack and the support
reserve of Section~\ref{sec:mixed-activity-prices}. Its tree
degree condition holds eventually because the geometric
neighbour count has upper exponent zero whereas the vertex
price has negative upper exponent. This supplies the
differentiated product input for the connected-support estimates.
It preserves the incoming representation and its selected
history; it does not construct or reset their parameters.
